% GENERATED FILE. Edit canonical lessons, then run npm run generate:books.
% canonical-source-hash: fnv1a64:66fbeb6b2ccd2ab2
% canonical-lessons: UR-C436-qudrat, UR-C436-khushbu, UR-C436-konpal, UR-C436-pankhuri, UR-C436-murjhana

\chapter{Nature}
\label{ch:ur-nature}
\hlchaptermodality{436}

\begin{chapteropening}
\textbf{By the end of this chapter:} \emph{I can talk about the natural world: scent, buds, petals and flowers that fade.}

Complete the last lesson of chapter 436: I can talk about the natural world: scent, buds, petals and flowers that fade.
\end{chapteropening}

\section[\texorpdfstring{\ur{قدرت} (qudrat)}{قدرت (qudrat)}]{\ur{قدرت} (qudrat) — nature; the natural world}

\label{lesson:UR-C436-qudrat}



\begin{quote}
Before the new one: say the Urdu for a seashell, then the Urdu for a dip (under water).
\end{quote}

\subsection*{You'll want to know: \ur{قدرت}}
\textbf{\ur{قدرت}} — \emph{qudrat} — \textquotedblleft{}nature; the natural world\textquotedblright{}.

\begin{quote}
\ur{یہ} \ur{وادی} \ur{قدرت} \ur{کا} \ur{تحفہ} \ur{ہے۔}
\end{quote}

\begin{quote}
\emph{yih vādī qudrat kā tohfā hai.}
\end{quote}

\begin{quote}
This valley is a gift of nature.
\end{quote}

Feminine.

\begin{cousinweb}
\textbf{\ur{قدرت}} (\emph{qudrat}) is usually traced, through Persian, to Arabic \emph{qudra}, \textquotedblleft{}power, might\textquotedblright{}; in Urdu it is also the natural world.
\end{cousinweb}

\begin{grammarlens}[title={What you've built}]
One more word for getting around.
\end{grammarlens}

\subsection*{Your turn}
\begin{itemize}
\raggedright
  \item \emph{Say it:} \emph{qudrat}
  \item \emph{Say it:} \emph{qudrat}, once more
  \item \emph{Say it:} \emph{yih vādī qudrat kā tohfā hai.} — This valley is a gift of nature.
  \item \emph{Recall:} say the Urdu for a seashell, then the Urdu for a dip (under water), then say \emph{qudrat} again
\end{itemize}

\begin{tcolorbox}[breakable,colback=teal!4,colframe=teal!35!black,title={Before you move on}]
What does \ur{قدرت} mean? (\textquotedblleft{}Nature; the natural world\textquotedblright{}.) Say it once more.
\end{tcolorbox}
\section[\texorpdfstring{\ur{خوشبو} (khushbū)}{خوشبو (khushbū)}]{\ur{خوشبو} (khushbū) — a fragrance; a scent}

\label{lesson:UR-C436-khushbu}



\begin{quote}
Before the new one: say the Urdu for a dip (under water), then the Urdu for nature.
\end{quote}

\subsection*{You'll want to know: \ur{خوشبو}}
\textbf{\ur{خوشبو}} — \emph{khushbū} — \textquotedblleft{}a fragrance; a scent\textquotedblright{}.

\begin{quote}
\ur{پھولوں} \ur{کی} \ur{خوشبو} \ur{کمرے} \ur{تک} \ur{آتی} \ur{ہے۔}
\end{quote}

\begin{quote}
\emph{phūloṅ kī khushbū kamre tak ātī hai.}
\end{quote}

\begin{quote}
The scent of the flowers reaches the room.
\end{quote}

Feminine.

\begin{cousinweb}
\textbf{\ur{خوشبو}} (\emph{khushbū}) is built from Persian \emph{khush}, \textquotedblleft{}pleasant\textquotedblright{}, and \emph{bū}, \textquotedblleft{}a smell\textquotedblright{}: it is usually traced to Persian \emph{khushbū}; \textbf{\ur{بدبو}} is its opposite.
\end{cousinweb}

\begin{grammarlens}[title={What you've built}]
One more word for getting around.
\end{grammarlens}

\subsection*{Your turn}
\begin{itemize}
\raggedright
  \item \emph{Say it:} \emph{khushbū}
  \item \emph{Say it:} \emph{khushbū}, once more
  \item \emph{Say it:} \emph{phūloṅ kī khushbū kamre tak ātī hai.} — The scent of the flowers reaches the room.
  \item \emph{Recall:} say the Urdu for a dip (under water), then the Urdu for nature, then say \emph{khushbū} again
\end{itemize}

\begin{tcolorbox}[breakable,colback=teal!4,colframe=teal!35!black,title={Before you move on}]
What does \ur{خوشبو} mean? (\textquotedblleft{}A fragrance; a scent\textquotedblright{}.) Say it once more.
\end{tcolorbox}
\section[\texorpdfstring{\ur{کونپل} (koṅpal)}{کونپل (koṅpal)}]{\ur{کونپل} (koṅpal) — a new shoot; an opening bud}

\label{lesson:UR-C436-konpal}



\begin{quote}
Before the new one: say the Urdu for nature, then the Urdu for a fragrance.
\end{quote}

\subsection*{You'll want to know: \ur{کونپل}}
\textbf{\ur{کونپل}} — \emph{koṅpal} — \textquotedblleft{}a new shoot; an opening bud\textquotedblright{}.

\begin{quote}
\ur{بہار} \ur{میں} \ur{درختوں} \ur{پر} \ur{کونپلیں} \ur{آتی} \ur{ہیں۔}
\end{quote}

\begin{quote}
\emph{bahār meṅ darakhtoṅ par koṅpaleṅ ātī haiṅ.}
\end{quote}

\begin{quote}
In spring new shoots appear on the trees.
\end{quote}

Feminine.

\begin{cousinweb}
\textbf{\ur{کونپل}} (\emph{koṅpal}) is inherited: CDIAL lists the Hindi cousin \emph{kõpal}, \textquotedblleft{}an opening bud\textquotedblright{}, under Sanskrit \emph{kuḍmalá}, \textquotedblleft{}a bud\textquotedblright{} (CDIAL 3250).
\end{cousinweb}

\begin{grammarlens}[title={What you've built}]
One more word for getting around.
\end{grammarlens}

\subsection*{Your turn}
\begin{itemize}
\raggedright
  \item \emph{Say it:} \emph{koṅpal}
  \item \emph{Say it:} \emph{koṅpal}, once more
  \item \emph{Say it:} \emph{bahār meṅ darakhtoṅ par koṅpaleṅ ātī haiṅ.} — In spring new shoots appear on the trees.
  \item \emph{Recall:} say the Urdu for nature, then the Urdu for a fragrance, then say \emph{koṅpal} again
\end{itemize}

\begin{tcolorbox}[breakable,colback=teal!4,colframe=teal!35!black,title={Before you move on}]
What does \ur{کونپل} mean? (\textquotedblleft{}A new shoot; an opening bud\textquotedblright{}.) Say it once more.
\end{tcolorbox}
\section[\texorpdfstring{\ur{پنکھڑی} (pankhuṛī)}{پنکھڑی (pankhuṛī)}]{\ur{پنکھڑی} (pankhuṛī) — a petal}

\label{lesson:UR-C436-pankhuri}



\begin{quote}
Before the new one: say the Urdu for a fragrance, then the Urdu for a new shoot.
\end{quote}

\subsection*{You'll want to know: \ur{پنکھڑی}}
\textbf{\ur{پنکھڑی}} — \emph{pankhuṛī} — \textquotedblleft{}a petal\textquotedblright{}.

\begin{quote}
\ur{اس} \ur{پھول} \ur{کی} \ur{پنکھڑیاں} \ur{نرم} \ur{اور} \ur{لال} \ur{ہیں۔}
\end{quote}

\begin{quote}
\emph{is phūl kī pankhuṛiyāṅ narm aur lāl haiṅ.}
\end{quote}

\begin{quote}
This flower's petals are soft and red.
\end{quote}

Feminine.

\begin{cousinweb}
\textbf{\ur{پنکھڑی}} (\emph{pankhuṛī}) is inherited: CDIAL lists the Hindi cousins \emph{pakhuṛī} and \emph{pā̃khṛī}, \textquotedblleft{}a petal\textquotedblright{}, under Sanskrit \emph{pákṣman}, \textquotedblleft{}eyelashes, the petal of a flower\textquotedblright{} (CDIAL 7638).
\end{cousinweb}

\begin{grammarlens}[title={What you've built}]
One more word for getting around.
\end{grammarlens}

\subsection*{Your turn}
\begin{itemize}
\raggedright
  \item \emph{Say it:} \emph{pankhuṛī}
  \item \emph{Say it:} \emph{pankhuṛī}, once more
  \item \emph{Say it:} \emph{is phūl kī pankhuṛiyāṅ narm aur lāl haiṅ.} — This flower's petals are soft and red.
  \item \emph{Recall:} say the Urdu for a fragrance, then the Urdu for a new shoot, then say \emph{pankhuṛī} again
\end{itemize}

\begin{tcolorbox}[breakable,colback=teal!4,colframe=teal!35!black,title={Before you move on}]
What does \ur{پنکھڑی} mean? (\textquotedblleft{}A petal\textquotedblright{}.) Say it once more.
\end{tcolorbox}
\section[\texorpdfstring{\ur{مرجھانا} (murjhānā)}{مرجھانا (murjhānā)}]{\ur{مرجھانا} (murjhānā) — to wilt; to fade}

\label{lesson:UR-C436-murjhana}



\begin{quote}
Before the new one: say the Urdu for a new shoot, then the Urdu for a petal.
\end{quote}

\subsection*{You'll want to know: \ur{مرجھانا}}
\textbf{\ur{مرجھانا}} — \emph{murjhānā} — \textquotedblleft{}to wilt; to fade\textquotedblright{}.

\begin{quote}
\ur{گرمی} \ur{میں} \ur{پودے} \ur{مرجھا} \ur{گئے۔}
\end{quote}

\begin{quote}
\emph{garmī meṅ paude murjhā gae.}
\end{quote}

\begin{quote}
The plants wilted in the heat.
\end{quote}

Past (it …): \ur{وہ} \ur{پھول} \ur{مرجھا} \ur{گیا}, \emph{voh phūl murjhā gayā}.

\begin{cousinweb}
\textbf{\ur{مرجھانا}} (\emph{murjhānā}) is an everyday word whose source the dictionaries we use do not give: CDIAL has no entry for it. No dictionary we use records an older form, so its history is not settled.
\end{cousinweb}

\begin{grammarlens}[title={What you've built}]
One more word for getting around.
\end{grammarlens}

\subsection*{Your turn}
\begin{itemize}
\raggedright
  \item \emph{Say it:} \emph{murjhānā}
  \item \emph{Say it:} \emph{murjhānā}, once more
  \item \emph{Say it:} \emph{garmī meṅ paude murjhā gae.} — The plants wilted in the heat.
  \item \emph{Recall:} say the Urdu for a new shoot, then the Urdu for a petal, then say \emph{murjhānā} again
\end{itemize}

\begin{tcolorbox}[breakable,colback=teal!4,colframe=teal!35!black,title={Before you move on}]
What does \ur{مرجھانا} mean? (\textquotedblleft{}To wilt; to fade\textquotedblright{}.) Say it once more.
\end{tcolorbox}
